\documentclass[11pt,a4paper]{article}

% Packages
\usepackage[utf8]{inputenc}
\usepackage[T1]{fontenc}
\usepackage[margin=0.75in]{geometry}
\usepackage{enumitem}
\usepackage{hyperref}
\usepackage{titlesec}
\usepackage{xcolor}

% Formatting
\pagestyle{empty}
\setlength{\parindent}{0pt}
\setlist[itemize]{leftmargin=*, noitemsep, topsep=3pt}

% Section formatting
\titleformat{\section}{\Large\bfseries}{}{0em}{}[\titlerule]
\titlespacing*{\section}{0pt}{12pt}{6pt}

\titleformat{\subsection}[runin]{\bfseries}{}{0em}{}
\titlespacing*{\subsection}{0pt}{8pt}{0pt}

% Colors
\definecolor{linkcolor}{RGB}{0,102,204}
\hypersetup{
    colorlinks=true,
    urlcolor=linkcolor
}

\begin{document}

% Header
\begin{center}
    {\LARGE \textbf{ERIC MAYEUX}}\\
    \vspace{4pt}
    {\large Chargé de projet sénior TI | Senior IT Project Manager}\\
    \vspace{4pt}
    Montréal, QC, Canada\\
    \href{mailto:mayeux.e@gmail.com}{mayeux.e@gmail.com} | 438-884-7645
\end{center}

\vspace{8pt}

\section*{Profil Professionnel}
\noindent
\textbf{Chargé de projet sénior TI} avec 10+ ans d'expérience en \textbf{gestion de projets complexes} dans le secteur financier (BNC). Expertise en \textbf{gouvernance de projets stratégiques}, coordination inter-équipes et \textbf{reddition de comptes}. Maîtrise des projets à fort contexte organisationnel : \textbf{gestion des risques}, livrables stratégiques, arrimage multi-partenaires. Certifié \textbf{PMP} (en cours), CSM. Agile-Waterfall hybride selon le contexte.

\section*{Compétences Techniques}

\textbf{Gouvernance \& Coordination:} Gouvernance de projets, planification détaillée, livrables stratégiques, coordination inter-partenaires, reddition de comptes, tableaux de bord décisionnels

\textbf{Gestion des risques:} Évaluation, documentation, plans de contingence, suivi continu, gestion proactive des enjeux, gestion des changements

\textbf{Parties prenantes:} Coordination multi-équipes (internes/externes), communication haute direction (VPTI), alignement sur les objectifs stratégiques

\textbf{Outils:} Jira, Confluence, MS Project, Datadog, Splunk, tableaux de bord KPI

\textbf{Méthodologies:} Agile, Scrum, SAFe, Kanban, Waterfall, hybride, PMP, PMI-ACP

\textbf{Intelligence Artificielle:} Agents IA, automatisation de tâches de gestion, RPA

\textbf{Architecture \& Cloud:} AWS Cloud Practitioner, Architectures Distribuées, intégration multi-systèmes

\vspace{6pt}

\section*{Réalisations}
\begin{itemize}
    \item Pilotage de projets TI stratégiques dans un contexte bancaire à haute \textbf{gouvernance} (BNC)
    \item Coordination de \textbf{multiples partenaires} internes et externes dans un contexte SAFe (Release Train)
    \item Mise en place de \textbf{rapports de suivi, bilans et tableaux de bord} pour la haute direction
    \item \textbf{Gestion proactive des risques} et enjeux : identification, documentation, actions correctives
    \item Supervision des \textbf{transferts en exploitation} et planification des transitions BAU
    \item Organisation de hackathons et POC : innovation, IA, migration, amélioration continue
\end{itemize}

\section*{Expérience Professionnelle}

\subsection*{Scrum Master -- Gestion de projet | Project Manager} \hfill \textit{Novembre 2025 -- Présent}\\
\textbf{Banque Nationale du Canada (BNC)} -- Montréal, QC\\
Pilote des projets TI à haute \textbf{gouvernance} dans un contexte organisationnel complexe.
\begin{itemize}
    \item Assure la \textbf{vision commune}, la gouvernance et la planification des projets auprès des parties prenantes
    \item Coordonne les équipes IT multidisciplinaires ; gère les \textbf{dépendances et interdépendances}
    \item \textbf{Gère les risques} et enjeux proactivement : plans de mitigation, communication adaptée à la haute direction
    \item Produit des \textbf{rapports de suivi, tableaux de bord} et dossiers décisionnels pour les instances de gouvernance
    \item Formule des \textbf{recommandations stratégiques} et pilote les ajustements de portée
    \item Développe des agents IA pour l'\textbf{automatisation} des tâches administratives et de gestion
\end{itemize}

\subsection*{Intégrateur Sénior Qualité | Senior Quality Integrator} \hfill \textit{Octobre 2023 -- Novembre 2025}\\
\textbf{Banque Nationale du Canada (BNC)} -- Montréal, QC
\begin{itemize}
    \item Coordination multi-sites (Thaïlande) : gestion des dépendances, \textbf{reddition de comptes}
    \item Tableaux de bord (Datadog, Splunk, Jira) et \textbf{métriques} pour la haute direction
    \item Pilotage de POC et projets d'innovation : Playwright MCP, Crew AI, hackathon
\end{itemize}

\subsection*{Intégrateur Qualité -- SDET | Quality Integrator -- Software Development Engineer in Test} \hfill \textit{Janvier 2022 -- Octobre 2023}\\
\textbf{Banque Nationale du Canada (BNC)} via Groupe Astek -- Montréal, QC
\begin{itemize}
    \item \textbf{Gestion des risques} qualité : analyse statique de sécurité, tests de charge, plans de contingence
    \item Stack : Selenium, AppliTools, RestAssured, JMeter, Gatling, K6
\end{itemize}

\subsection*{QA Automaticien | QA Automation Engineer} \hfill \textit{Juin 2021 -- Décembre 2021}\\
\textbf{AODocs} -- Paris, France
\begin{itemize}
    \item Automatisation tests non-régression ; POC Flutter ; rapports Allure
\end{itemize}

\subsection*{Automaticien CI/CD | DevOps Automation Engineer} \hfill \textit{Janvier 2020 -- Juin 2021}\\
\textbf{ENGIE DIGITAL} via Groupe Hénix -- Paris, France
\begin{itemize}
    \item Création CI/CD : Jenkins, Docker ; administration Jira, KPIs ; scripting Python/Shell
\end{itemize}

\subsection*{Product Owner} \hfill \textit{Juin 2019 -- Septembre 2019}\\
\textbf{Hénix} -- Montrouge, France
\begin{itemize}
    \item Gestion backlog, priorisation, support clients SaaS, coordination partenaires Jira/Squash
\end{itemize}

\subsection*{Développeur Logiciel | Software Developer} \hfill \textit{Mai 2017 -- Mai 2019}\\
\textbf{MEDIAPOST} via Groupe Hénix -- Paris, France
\begin{itemize}
    \item Développement back office ; intégration multi-systèmes (SOAP/REST, PL/SQL, Sybase)
\end{itemize}

\subsection*{Consultant QA \& Automatisation | QA Automation Consultant} \hfill \textit{Avril 2016 -- Avril 2017}\\
\textbf{ENGIE IT} via Groupe Hénix -- Saint-Ouen, France
\begin{itemize}
    \item Gestion opérationnelle des environnements ; coordination prestataires et partenaires
\end{itemize}

\section*{Certifications \& Formations}

\textbf{PMP Certification} (en cours) \hfill \textit{2026}

Project Management Professional -- Project Management Institute (PMI)

\textbf{AWS Cloud Practitioner} \hfill \textit{2026}

Amazon Web Services (AWS)

\textbf{Certified Scrum Master (CSM)} \hfill \textit{2024}

Scrum Alliance

\textbf{Jira ACP-600 -- Project Administration in Jira Server} \hfill \textit{2019}

Atlassian Certified Professional

\textbf{Cursus Automaticien et DevOps} \hfill \textit{2019}\\
\textit{AFCEPF -- Montrouge, France}

\textbf{Cursus Architecture logiciel \& Analyste informaticien - Certifié Master II} \hfill \textit{2017}\\
\textit{AFCEPF-HENIX -- Malakoff, France}

\section*{Formation Académique}

\textbf{Licence en Gestion (Bachelor's Degree in Management)} \hfill \textit{2011}\\
École Supérieure de Gestion (E.S.G) -- Paris, France

\section*{Compétences Linguistiques}

\textbf{Français:} Langue maternelle (Native)\\
\textbf{Anglais:} Niveau Avancé (Advanced / Professional Working Proficiency)

\end{document}
